% titel: Blatt 10, A2: Säurestärke-Reihen
% skript: Mortimer Kap. 17 (Säurestärke)

% Teil I: Werkzeuge, kein Lösungsweg.

Die Säurestärke (Häufigkeit der H$^+$-Abgabe) hängt von der Stabilität der
entstehenden Base und von der Polarität der X--H-Bindung ab.

\begin{werkzeugkasten}[Säurestärke abschätzen]
\begin{itemize}
  \item \textbf{Halogenwasserstoffe HX:} Stärke wächst mit der Grösse/Abnahme
        der H--X-Bindungsenergie, also HF $<$ HCl $<$ HBr $<$ HI.
  \item \textbf{Metallhydride wie NaH} sind sehr starke Basen, keine Säuren
        (H$^-$-Donator).
  \item \textbf{Sauerstoffsäuren H$_x$YO$_m$:} je mehr Sauerstoffatome in den
        Zentralteil eingebunden sind (höhere O-Zahl pro Zentralatom), desto
        stärker die Säure (destabilisiert die Bindung, stabilisiert das
        Anion). Beispiele: H$_2$SeO$_4$ $>$ H$_2$SeO$_3$.
  \item Neutrales H$_2$O ist eine sehr schwache Säure.
  \item Elektronegativere Zentralatome mit höherer Oxidationsstufe machen
        die Säure stärker.
\end{itemize}
\textit{Skript:} Mortimer Kap. 17 (Säurestärke).
\end{werkzeugkasten}

\begin{achtung}[Darauf achten]
\begin{itemize}
  \item NaH ist ein Salz mit H$^-$ --- keine Brønsted-Säure; in der Reihe
        betrachtet man es als starke Base/keine Säure.
  \item HF ist unter den Halogenwasserstoffen die \emph{schwächste} (weil
        die H-F-Bindung am stabilsten), nicht die stärkste.
  \item H$_2$SeO$_4$ (mehr O) ist stärker als H$_2$SeO$_3$.
  \item H$_2$O steht als sehr schwache Säure am Ende.
\end{itemize}
\end{achtung}
